\subsection{Justificación del proyecto de innovación / mejora / Creatividad}

El desarrollo e implementación del sistema web con flujos automatizados para la gestión de nóminas y recursos humanos en la empresa J\&P Periféricos S.A.C. 
responde a la necesidad de sustituir el procesamiento manual de haberes por una plataforma robusta, auditable y alineada a las exigencias normativas del sector productivo. 
El alcance del proyecto se enfoca en los módulos críticos de la organización (autenticación, gestión de legajos, asistencia/incidencias, motor de cálculo salarial, dispersión contable y auditoría), postergando el subsistema de autoservicio del colaborador como una línea de desarrollo futuro. A continuación, se detallan las justificaciones que sustentan la viabilidad y valor de la propuesta:

\subsubsection{Justificación Práctica}
La utilidad práctica del proyecto consiste en integrar la información administrada por los responsables de Recursos Humanos con las operaciones de revisión y cálculo realizadas por Contabilidad. 
La solución permitirá gestionar legajos e incidencias, validar datos, revisar resultados y generar documentos dentro de un proceso con responsabilidades diferenciadas. 
Para J\&P, el desarrollo constituye una base tecnológica para ampliar su portafolio de soluciones empresariales.

\subsubsection{Justificación Tecnológica}
El proyecto se justifica tecnológicamente por el desarrollo de una plataforma web que integra la información laboral, 
el procesamiento de nóminas y los controles administrativos en una solución que J\&P Periféricos S.A.C. proyecta ofrecer por suscripción a clientes empresariales. 
Su arquitectura de monolito modular organiza las funcionalidades de Recursos Humanos, Contabilidad y auditoría en módulos con responsabilidades diferenciadas, facilitando el mantenimiento y la incorporación progresiva de nuevas funciones.

El uso de TypeScript y NestJS permite estructurar la lógica del sistema e implementar validaciones y controles de acceso según las responsabilidades de cada usuario. PostgreSQL, junto con
Prisma ORM, permite gestionar información relacionada y ejecutar operaciones transaccionales para mantener la consistencia de los registros. La exactitud de los cálculos remunerativos se comprobará
mediante pruebas con resultados esperados, mientras que el registrode modificaciones permitirá identificar quién realizó cada cambio, cuándo se efectuó y qué datos fueron modificados.

Asimismo, n8n permitirá coordinar flujos automatizados para la preparación de archivos de pago y la distribución de boletas digitales. 
Redis apoyará el procesamiento de tareas en segundo plano, y Docker facilitará el empaquetado y despliegue de los servicios. 
Se propone alojar la solución en infraestructura de nube, con una configuración que será definida y validada durante la implementación. 
Este enfoque constituye una base técnica para la futura prestación del servicio por suscripción, cuya operación requerirá mantenimiento, respaldo y supervisión de los servicios.

\subsubsection{Justificación Legal}
El desarrollo del sistema está justificado en el plano legal, dado que su funcionamiento permite la automatización de los procesos de cálculo de planillas y, a su vez, asegura el seguimiento del cumplimiento del marco regulatorio del sector laboral privado en el Perú.

Decreto Supremo N.º 054-97-EF: Dictamina la afiliación obligatoria de los empleados al sistema de pensiones (ONP/AFP), además de designar al empleador como agente de retención de los aportes a los regímenes pensionarios.

Decreto Supremo N.º 179-2004-EF: Delega la obligación a la empresa de ser el agente de retención, descontando el Impuesto a la Renta de 5.ª Categoría de los empleados que apliquen a este tramo bajo una escala progresiva.

Decreto Supremo N.º 018-2007-TR: Determina el uso obligatorio de la planilla electrónica.

Ley N.º 29903: Afiliación obligatoria de nuevos trabajadores al régimen de AFP.

Ley N.º 26790: Afiliación obligatoria de los trabajadores a EsSalud, con la aportación del 9\% asumida por el empleador.

Ley N.º 27735: Derecho del trabajador a percibir dos gratificaciones al año (julio y diciembre), equivalentes a un sueldo completo cada una.

Ley N.º 25129: Dictamina el pago de una remuneración extra del 10\% de la RMV (Remuneración Mínima Vital) a los trabajadores que tengan hijos menores de edad o dependientes que cursen estudios superiores

\subsubsection{Justificación Económica}
Para J\&P Periféricos S.A.C., el proyecto representa una oportunidad de diversificar su portafolio mediante una solución que se proyecta comercializar por suscripción. 
La evaluación económica considerará los costos de desarrollo, alojamiento, mantenimiento y soporte, frente a escenarios estimados de ingresos. Los precios, la cantidad de suscriptores y los beneficios económicos se identificarán como supuestos de proyección mientras no existan contratos o resultados comerciales verificables.

\subsubsection{Justificación Personal}
A nivel formativo y profesional, este proyecto de titulación representa la oportunidad de consolidar, aplicar y validar las capacidades tecnicas, arquitectonicas y de gestión adquiridos a lo largo de la formación académica y las prácticas empresarial. 
El desarrollo integral de una solución empresarial desde cero asumiendo la investigación de los casos de uso, las tecnologias a utilizar y el diseño de la base de datos constituye a un reto tecnico que fortalecera el perfil profesional de los autores como desarrolladores de software.